% Options for packages loaded elsewhere
\PassOptionsToPackage{unicode}{hyperref}
\PassOptionsToPackage{hyphens}{url}
\documentclass[
  ignorenonframetext,
]{beamer}
\newif\ifbibliography
\usepackage{pgfpages}
\setbeamertemplate{caption}[numbered]
\setbeamertemplate{caption label separator}{: }
\setbeamercolor{caption name}{fg=normal text.fg}
\beamertemplatenavigationsymbolsempty
% remove section numbering
\setbeamertemplate{part page}{
  \centering
  \begin{beamercolorbox}[sep=16pt,center]{part title}
    \usebeamerfont{part title}\insertpart\par
  \end{beamercolorbox}
}
\setbeamertemplate{section page}{
  \centering
  \begin{beamercolorbox}[sep=12pt,center]{section title}
    \usebeamerfont{section title}\insertsection\par
  \end{beamercolorbox}
}
\setbeamertemplate{subsection page}{
  \centering
  \begin{beamercolorbox}[sep=8pt,center]{subsection title}
    \usebeamerfont{subsection title}\insertsubsection\par
  \end{beamercolorbox}
}
% Prevent slide breaks in the middle of a paragraph
\widowpenalties 1 10000
\raggedbottom
\AtBeginPart{
  \frame{\partpage}
}
\AtBeginSection{
  \ifbibliography
  \else
    \frame{\sectionpage}
  \fi
}
\AtBeginSubsection{
  \frame{\subsectionpage}
}
\usepackage{iftex}
\ifPDFTeX
  \usepackage[T1]{fontenc}
  \usepackage[utf8]{inputenc}
  \usepackage{textcomp} % provide euro and other symbols
\else % if luatex or xetex
  \usepackage{unicode-math} % this also loads fontspec
  \defaultfontfeatures{Scale=MatchLowercase}
  \defaultfontfeatures[\rmfamily]{Ligatures=TeX,Scale=1}
\fi
\usepackage{lmodern}
\ifPDFTeX\else
  % xetex/luatex font selection
\fi
% Use upquote if available, for straight quotes in verbatim environments
\IfFileExists{upquote.sty}{\usepackage{upquote}}{}
\IfFileExists{microtype.sty}{% use microtype if available
  \usepackage[]{microtype}
  \UseMicrotypeSet[protrusion]{basicmath} % disable protrusion for tt fonts
}{}
\makeatletter
\@ifundefined{KOMAClassName}{% if non-KOMA class
  \IfFileExists{parskip.sty}{%
    \usepackage{parskip}
  }{% else
    \setlength{\parindent}{0pt}
    \setlength{\parskip}{6pt plus 2pt minus 1pt}}
}{% if KOMA class
  \KOMAoptions{parskip=half}}
\makeatother
\setlength{\emergencystretch}{3em} % prevent overfull lines
\providecommand{\tightlist}{%
  \setlength{\itemsep}{0pt}\setlength{\parskip}{0pt}}
% Auto-generated by infrastructure/rendering/_pdf_latex_helpers.py::extract_math_font_preamble.
% Minimal header passed to Pandoc via -H so Beamer slides pick up the same math font
% as the combined-PDF path. Adjust the manuscript's preamble.md to change the font globally.
\usepackage{unicode-math}
\setmathfont{latinmodern-math.otf}

% Slide layout defaults for warning-clean scientific decks.
\usepackage{etoolbox}
\IfFileExists{xurl.sty}{\usepackage{xurl}}{}
\IfFileExists{seqsplit.sty}{\usepackage{seqsplit}}{\newcommand{\seqsplit}[1]{#1}}
\protected\def\breakseq#1{\seqsplit{#1}}
\protected\def\breaktt#1{\begingroup\ttfamily\seqsplit{#1}\endgroup}
\setlength{\emergencystretch}{6em}
\tolerance=5000
\hbadness=10000
\hfuzz=1pt
\setlength{\tabcolsep}{2pt}
\AtBeginEnvironment{longtable}{\tiny\renewcommand{\arraystretch}{0.86}\setlength{\tabcolsep}{1pt}}
\AtBeginEnvironment{tabular}{\tiny\renewcommand{\arraystretch}{0.86}\setlength{\tabcolsep}{1pt}}

% Natbib and cross-reference fallbacks — slides don't load natbib
% or cleveref, but combined-PDF manuscript prose may emit these
% commands. The fallback renders citations as a bracketed key list
% and cross-references as detokenized labels so slides don't fail on
% undefined control sequences or raw underscores. \providecommand is
% a no-op if packages load later.
\providecommand{\citep}[1]{[#1]}
\providecommand{\citet}[1]{#1}
\providecommand{\citealp}[1]{#1}
\providecommand{\citeauthor}[1]{#1}
\providecommand{\citeyear}[1]{#1}
\providecommand{\cref}[1]{\texttt{\detokenize{#1}}}
\providecommand{\Cref}[1]{\texttt{\detokenize{#1}}}
\usepackage{bookmark}
\IfFileExists{xurl.sty}{\usepackage{xurl}}{} % add URL line breaks if available
\urlstyle{same}
% fallback for those not using the hyperref driver hyperxmp:
\makeatletter
\@ifundefined{xmpquote}{\newcommand{\xmpquote}[1]{#1}}{}
\makeatother
\hypersetup{
  hidelinks,
  pdfcreator={LaTeX via pandoc}}

\author{\texorpdfstring{}{}}
\date{}

\begin{document}

\section{Introduction}\label{sec:introduction}

\begin{frame}[fragile,allowframebreaks]{Introduction}
This \breaktt{template\_code\_project} serves as the foundational
exemplar for the \href{https://github.com/docxology/template}{Research
Project Template} ecosystem, demonstrating a fully-tested numerical
optimization implementation bracketed by rigorous infrastructure,
hermetic testing, and extensive documentation architectures. The prose,
the labelled figures, and the labelled equations have all been generated
through an auditable custody chain starting from algorithm
implementation through strict CI/CD validation to multi-format
\texttt{.pdf} compilation.
\end{frame}

\begin{frame}[fragile,allowframebreaks]{Template Architecture Context}
\protect\phantomsection\label{template-architecture-context}
Scientific engineering requires mathematical accuracy combined with
software reliability. This project unifies theoretical optimization with
the repository's three foundational pillars:

\begin{enumerate}
\tightlist
\item
  \textbf{\texttt{infrastructure/} Layer (Root Directory)}: A modular
  stack of importable Python packages providing the computational
  scaffolding. The current package count is measured in the template
  repository's generated canonical facts rather than repeated here
  because it changes as infrastructure modules are added or retired.
\item
  \textbf{\texttt{tests/} Framework
  (\breaktt{projects/templates/template\_code\_project/tests/})}: An
  uncompromising validation layer maintaining a zero-mock testing
  policy. This is enforced automatically via the
  \href{https://github.com/docxology/template/blob/main/.github/workflows/ci.yml}{CI
  workflow} mapping to \texttt{pyproject.toml} directives.
\item
  \textbf{\texttt{docs/} Knowledge Base
  (\breaktt{projects/templates/template\_code\_project/docs/})}: A
  structured repository of architectural guidelines, operational
  patterns, and the Rigorous Agentic Scientific Protocol (RASP) that
  governs the AI-assisted agents writing these very texts.
\end{enumerate}

This implementation of gradient descent algorithms for solving
optimization problems is used as the vehicle to demonstrate these
pillars. The theoretical problem stated in
\breakseq{[@eq:optimization\_problem]} is mapped programmatically inside the
\href{https://github.com/docxology/template/blob/main/projects/templates/template_code_project/src/optimizer.py}{optimizer
module}:

\begin{equation}
\label{eq:optimization_problem}
\min_{x \in \mathbb{R}^n} f(x)
\end{equation}

where \(f: \mathbb{R}^n \rightarrow \mathbb{R}\) is a continuously
differentiable objective function.
\end{frame}

\begin{frame}[fragile,allowframebreaks]{Infrastructure Integration}
\protect\phantomsection\label{infrastructure-integration}
Rather than existing as isolated scripts, this project extensively
leverages the \texttt{infrastructure} layer:

\begin{itemize}
\tightlist
\item
  \textbf{Scientific Utilities}: Utilizing
  \breaktt{infrastructure.scientific.stability} and
  \breaktt{infrastructure.scientific.benchmarking} to guarantee numerical
  boundaries and performance scaling.
\item
  \textbf{Hermetic Validation}: Deploying
  \breaktt{infrastructure.validation} components
  (\breaktt{markdown\_validator}, \breaktt{output\_validator}) to ensure
  generated artifacts are structurally valid and traceable.
\item
  \textbf{Reporting \& Rendering}: Employing
  \breaktt{infrastructure.rendering.pdf\_renderer} and
  \breaktt{infrastructure.reporting.executive\_reporter} to automatically
  transform code outputs into this finalized manuscript.
\end{itemize}
\end{frame}

\begin{frame}[allowframebreaks]{Algorithm Overview}
\protect\phantomsection\label{algorithm-overview}
The reference gradient descent algorithm iteratively updates the
solution using the rule shown in \breakseq{[@eq:gradient\_descent\_update]}:

\begin{equation}
\label{eq:gradient_descent_update}
x_{k+1} = x_k - \alpha \nabla f(x_k)
\end{equation}

where \(\alpha > 0\) is the step size (learning rate) and
\(\nabla f(x_k)\) is the gradient of the objective function at iteration
\(k\).
\end{frame}

\begin{frame}[fragile,allowframebreaks]{Exemplar Implementation Goals}
\protect\phantomsection\label{exemplar-implementation-goals}
As the representative project for the repository, this implementation
explicitly demonstrates:

\begin{enumerate}
\tightlist
\item
  \textbf{Infrastructure-Coupled Code}: Scientific implementations that
  delegate logging, file ops, and reporting to the
  \texttt{infrastructure} core.
\item
  \textbf{Zero-Mock Verification}: A strict comprehensive validation
  suite proving numerical accuracy without artificial test boundaries.
\item
  \textbf{Automated Research Pipelines}: High-precision analyses that
  generate publication-quality, accessible visualizations automatically.
\item
  \textbf{Agentic Documentation standards}: Native adherence to the RASP
  methodology and \texttt{AGENTS.md} guidelines, ensuring the logic
  remains verifiable by both human and artificial intelligence.
\end{enumerate}
\end{frame}

\begin{frame}[fragile,allowframebreaks]{Reader's guide to the
manuscript}
\protect\phantomsection\label{readers-guide-to-the-manuscript}
\begin{itemize}
\tightlist
\item
  \textbf{\breakseq{[@sec:methodology]}} ties pseudocode to
  \breaktt{gradient\_descent()} and explains how stability checks and
  benchmarks call into \breaktt{infrastructure.scientific}.
\item
  \textbf{\breakseq{[@sec:results]}} is figure-centric: every panel references
  a generator in \texttt{src/figures/} (orchestrated via
  \breaktt{scripts/optimization\_analysis.py}) and uses
  \texttt{\{\{CONFIG\_*\}\}} / \texttt{\{\{RESULT\_*\}\}} placeholders
  for numeric values.
\item
  \textbf{\breakseq{[@sec:experimental\_setup]}} lists the exact YAML fields
  (\texttt{experiment:} block) that controlled the run whose artifacts
  you are viewing.
\item
  \textbf{\breakseq{[@sec:reproducibility]}} records the configuration hash and
  artifact inventory produced alongside the PDF.
\item
  \textbf{\breakseq{[@sec:scope]}} states scope and related literature so the
  exemplar is not mistaken for a general-purpose optimizer benchmark
  suite.
\end{itemize}
\end{frame}

\begin{frame}[allowframebreaks]{Why a quadratic model}
\protect\phantomsection\label{why-a-quadratic-model}
Restricting \(f\) to a quadratic with known \((A,b)\) keeps the optimum,
gradient, and spectral data explicit. For \(A = I\) and
\(b = \mathbf{1}\), the optimal point is \(x^\ast = \mathbf{1}\) and
gradient descent with fixed \(\alpha\) reduces to a linear iteration in
the error (see \breakseq{[@eq:scalar\_linear\_update]} in the results section).
That simplicity isolates step-size effects and makes divergent choices
(\(\alpha \geq 2\) in 1D) visible in both plots and CSV rows without
requiring trust-region or line-search fixes---those extensions are left
to future work (\breakseq{[@sec:scope]}).
\end{frame}

\end{document}
